% GENERATED FILE. Edit canonical lessons, then run npm run generate:books.
% canonical-source-hash: fnv1a64:c98adb70690e18d5
% canonical-lessons: MR-C198-dabne, MR-C198-khokne, MR-C198-thartharne, MR-C198-mat-dene, MR-C198-svagat-karne

\chapter{To press, To cough, To tremble, To vote, To welcome}
\label{ch:mr-to-press-to-cough-to-tremble-to-vote-to-welcome}
\hlchaptermodality{198}

\begin{chapteropening}
\textbf{By the end of this chapter:} \emph{I can say to press, to cough, to tremble, to vote and to welcome.}

Complete the last lesson of chapter 198: I can say to press, to cough, to tremble, to vote and to welcome.
\end{chapteropening}

\section[\texorpdfstring{\mr{दाबणे} (dābṇe)}{दाबणे (dābṇe)}]{\mr{दाबणे} (dābṇe) — to press}

\label{lesson:MR-C198-dabne}



\begin{quote}
Before the new one: say the Marathi for to chat, then the Marathi for to deliver.
\end{quote}

\subsection*{You'll want to know: \mr{दाबणे}}
\textbf{\mr{दाबणे}} — \emph{dābṇe} — \textquotedblleft{}to press\textquotedblright{}.

\begin{cousinweb}
\textbf{\mr{दाबणे}} (\emph{dābṇe}) comes from a reconstructed \emph{*dabb-}, \textquotedblleft{}press\textquotedblright{} (CDIAL 6173). CDIAL reports A. Master's view that it comes from Dravidian.
\end{cousinweb}

\begin{grammarlens}[title={What you've built}]
One more word for this chapter.
\end{grammarlens}

\subsection*{Your turn}
\begin{itemize}
\raggedright
  \item \emph{Say it:} \emph{dābṇe}
  \item \emph{Say it:} \emph{dābṇe}, once more
  \item \emph{Say it:} say \emph{dābṇe}
  \item \emph{Recall:} say the Marathi for to chat, then the Marathi for to deliver, then say \emph{dābṇe} again
\end{itemize}

\begin{tcolorbox}[breakable,colback=teal!4,colframe=teal!35!black,title={Before you move on}]
What does \mr{दाबणे} mean? (\textquotedblleft{}To press\textquotedblright{}.) Say it once more.
\end{tcolorbox}
\section[\texorpdfstring{\mr{खोकणे} (khokṇe)}{खोकणे (khokṇe)}]{\mr{खोकणे} (khokṇe) — to cough}

\label{lesson:MR-C198-khokne}



\begin{quote}
Before the new one: say the Marathi for to deliver, then the Marathi for to press.
\end{quote}

\subsection*{You'll want to know: \mr{खोकणे}}
\textbf{\mr{खोकणे}} — \emph{khokṇe} — \textquotedblleft{}to cough\textquotedblright{}.

\begin{cousinweb}
\textbf{\mr{खोकणे}} (\emph{khokṇe}) is a tadbhava, a word inherited and reshaped by sound change. It comes from \emph{*khokkh}, \textquotedblleft{}bark, cough\textquotedblright{}, a form CDIAL reconstructs from its descendants (CDIAL 3926).
\end{cousinweb}

\begin{grammarlens}[title={What you've built}]
One more word for this chapter.
\end{grammarlens}

\subsection*{Your turn}
\begin{itemize}
\raggedright
  \item \emph{Say it:} \emph{khokṇe}
  \item \emph{Say it:} \emph{khokṇe}, once more
  \item \emph{Say it:} say \emph{khokṇe}
  \item \emph{Recall:} say the Marathi for to deliver, then the Marathi for to press, then say \emph{khokṇe} again
\end{itemize}

\begin{tcolorbox}[breakable,colback=teal!4,colframe=teal!35!black,title={Before you move on}]
What does \mr{खोकणे} mean? (\textquotedblleft{}To cough\textquotedblright{}.) Say it once more.
\end{tcolorbox}
\section[\texorpdfstring{\mr{थरथरणे} (thartharṇe)}{थरथरणे (thartharṇe)}]{\mr{थरथरणे} (thartharṇe) — to tremble}

\label{lesson:MR-C198-thartharne}



\begin{quote}
Before the new one: say the Marathi for to press, then the Marathi for to cough.
\end{quote}

\subsection*{You'll want to know: \mr{थरथरणे}}
\textbf{\mr{थरथरणे}} — \emph{thartharṇe} — \textquotedblleft{}to tremble\textquotedblright{}.

\begin{cousinweb}
\textbf{\mr{थरथरणे}} (\emph{thartharṇe}) is a tadbhava, a word inherited and reshaped by sound change. It comes from \emph{*thar}, \textquotedblleft{}tremble\textquotedblright{}, a form CDIAL reconstructs from its descendants (CDIAL 6092). CDIAL compares Buddhist Sanskrit \emph{tharatharāyamāṇa}, \textquotedblleft{}trembling\textquotedblright{}.
\end{cousinweb}

\begin{grammarlens}[title={What you've built}]
One more word for this chapter.
\end{grammarlens}

\subsection*{Your turn}
\begin{itemize}
\raggedright
  \item \emph{Say it:} \emph{thartharṇe}
  \item \emph{Say it:} \emph{thartharṇe}, once more
  \item \emph{Say it:} say \emph{thartharṇe}
  \item \emph{Recall:} say the Marathi for to press, then the Marathi for to cough, then say \emph{thartharṇe} again
\end{itemize}

\begin{tcolorbox}[breakable,colback=teal!4,colframe=teal!35!black,title={Before you move on}]
What does \mr{थरथरणे} mean? (\textquotedblleft{}To tremble\textquotedblright{}.) Say it once more.
\end{tcolorbox}
\section[\texorpdfstring{\mr{मत} \mr{देणे} (mat deṇe)}{मत देणे (mat deṇe)}]{\mr{मत} \mr{देणे} (mat deṇe) — to vote}

\label{lesson:MR-C198-mat-dene}



\begin{quote}
Before the new one: say the Marathi for to cough, then the Marathi for to tremble.
\end{quote}

\subsection*{You'll want to know: \mr{मत} \mr{देणे}}
\textbf{\mr{मत} \mr{देणे}} — \emph{mat deṇe} — \textquotedblleft{}to vote\textquotedblright{}.

\begin{cousinweb}
\textbf{\mr{मत}} is Sanskrit \emph{mata}, \textquotedblleft{}thought, opinion\textquotedblright{}, a tatsama. \textbf{\mr{देणे}} comes from Sanskrit \emph{dadāti}, \textquotedblleft{}gives\textquotedblright{} (CDIAL 6141).
\end{cousinweb}

\begin{grammarlens}[title={What you've built}]
One more word for this chapter.
\end{grammarlens}

\subsection*{Your turn}
\begin{itemize}
\raggedright
  \item \emph{Say it:} \emph{mat deṇe}
  \item \emph{Say it:} \emph{mat deṇe}, once more
  \item \emph{Say it:} say \emph{mat deṇe}
  \item \emph{Recall:} say the Marathi for to cough, then the Marathi for to tremble, then say \emph{mat deṇe} again
\end{itemize}

\begin{tcolorbox}[breakable,colback=teal!4,colframe=teal!35!black,title={Before you move on}]
What does \mr{मत} \mr{देणे} mean? (\textquotedblleft{}To vote\textquotedblright{}.) Say it once more.
\end{tcolorbox}
\section[\texorpdfstring{\mr{स्वागत} \mr{करणे} (svāgat karṇe)}{स्वागत करणे (svāgat karṇe)}]{\mr{स्वागत} \mr{करणे} (svāgat karṇe) — to welcome}

\label{lesson:MR-C198-svagat-karne}



\begin{quote}
Before the new one: say the Marathi for to tremble, then the Marathi for to vote.
\end{quote}

\subsection*{You'll want to know: \mr{स्वागत} \mr{करणे}}
\textbf{\mr{स्वागत} \mr{करणे}} — \emph{svāgat karṇe} — \textquotedblleft{}to welcome\textquotedblright{}.

\begin{cousinweb}
\textbf{\mr{स्वागत}} is Sanskrit \emph{svāgata}, \textquotedblleft{}welcome\textquotedblright{}, a tatsama: \emph{su-}, \textquotedblleft{}well\textquotedblright{}, and \emph{āgata}, \textquotedblleft{}come\textquotedblright{}. \textbf{\mr{करणे}} is \textquotedblleft{}to do\textquotedblright{}.
\end{cousinweb}

\begin{grammarlens}[title={What you've built}]
One more word for this chapter.
\end{grammarlens}

\subsection*{Your turn}
\begin{itemize}
\raggedright
  \item \emph{Say it:} \emph{svāgat karṇe}
  \item \emph{Say it:} \emph{svāgat karṇe}, once more
  \item \emph{Say it:} say \emph{svāgat karṇe}
  \item \emph{Recall:} say the Marathi for to tremble, then the Marathi for to vote, then say \emph{svāgat karṇe} again
\end{itemize}

\begin{tcolorbox}[breakable,colback=teal!4,colframe=teal!35!black,title={Before you move on}]
What does \mr{स्वागत} \mr{करणे} mean? (\textquotedblleft{}To welcome\textquotedblright{}.) Say it once more.
\end{tcolorbox}
